\chapter{The Comparison}
% OUTLINE Ch5 "The Comparison (connections)". Laws: CITATION (Christoffel
% 1869; Levi-Civita 1917; Cartan 1923; Koszul via the standard texts),
% LENS (zero device appearances), Vol-1 laws (step-not-move watched).
% The rule as genuinely EXTRA structure (no treaty clause supplies it);
% the symbols honestly NOT a tensor; the Ch7 fundamental-theorem purchase
% seeded in one forward clause; the loop teaser hands Ch6 its subject.
% Gate: Kodo per-claim citation.

Two laboratories sit in neighboring villages, and each holds an arrow ---
a measured velocity, say, each valid at its own bench. A question that
sounds innocent ends the innocence of this book's construction: are the
two labs holding \emph{the same} arrow? On the flat paper of a schoolroom
the question never even registers as a question --- slide one arrow over
to the other, keeping it parallel to itself, and compare; everyone has
done it; nobody has ever been asked to justify it. But the sliding is the
justification, and on ground that is not flat it fails to be innocent:
carry an arrow along the surface of the Earth from the equator to the
pole by one route, and by another route, keeping it as parallel as the
ground permits at every stage, and the two deliveries arrive pointing in
different directions. The schoolroom's free slide was a hidden privilege
of flatness. On a general manifold, arrows at different points live in
different tangent spaces --- different rooms, as the last chapter said ---
and there is no built-in fact of the matter about which arrow here is
``the same'' as which arrow there. If the two labs want to compare, a
rule for carrying must be \emph{chosen}. This chapter is about that rule:
what choosing one means, what it costs, and who wrote its grammar.

State the honest position first, because the entire chapter rests on it:
nothing constructed so far supplies the rule. The treaty of Chapter 2
lets observers translate descriptions of one point; it says nothing about
relating one point to the next. The machines of Chapter 4 operate at
their stations. A rule for carrying arrows along curves --- the trade
calls it a connection, and the carrying itself parallel transport --- is
genuinely additional structure, laid on top of the smooth manifold by
decision, exactly as the topology was a declared inventory of regions and
the smooth structure a declared treaty of observers. The subject's
honesty lives in remembering which structures were declared and when.

What must a carrying rule provide? Along any curve, a prescription taking
an arrow at the starting point to an arrow at each later point, linearly
--- sums carry to sums, scalings to scalings --- and consistently with
the calculus already in hand. The working form of the demand is a
derivative: given a field of arrows and a direction to probe in, the rule
must produce the honest rate of change of the field --- the covariant
derivative --- obeying linearity and the Leibniz product rule. In a
chart, the naive gesture --- differentiate each component function and
call the results the derivative's components --- fails the observers'
treaty: transition maps bend, and the naive components of change do not
transform as an arrow's components must. The repair is a correction term,
one array of coefficients that records, chart by chart, how much the
rule's notion of ``held parallel'' disagrees with the chart's notion of
``components held constant.'' Those coefficients are the Christoffel
symbols, published by Elwin Bruno Christoffel in 1869 in his study of how
differential expressions transform --- and an honesty note belongs to
them immediately, because the trade's notation invites the error: the
symbols are not the components of any tensor. They are the coordinates of
the \emph{rule}, not readings of a machine; under a change of chart they
transform with an extra, inhomogeneous term --- the signature of a
dictionary, not of an object. Confusing the rule's bookkeeping with a
physical machine is the classic beginner's injury, and the cure is the
distinction this book has enforced since Chapter 3: machines transform by
the two laws; the symbols do not, because they are not machines.

The rule's geometric soul was exhibited by Tullio Levi-Civita in 1917,
and his construction deserves its picture. For a surface sitting inside
ordinary space, carry an arrow along a curve by the only procedure
available to an honest inhabitant: at each instant, let the arrow ride
unchanged in the surrounding space, then project away the component the
surface cannot hold --- keep what is tangent, discard what points out of
the world. The result --- parallel displacement, in his phrase --- is a
carrying rule derived from the surface's shape alone, and it is the rule
under which the equator-and-pole experiment above gives its
route-dependent answers. Levi-Civita's picture did for connections what
the diagrams of Chapter 4 did for tensors: it made the abstraction
walkable. The general concept --- a connection as a structure in its own
right, on any manifold, tied to no surrounding space and no metric ---
was forged by \'Elie Cartan in the early 1920s, his 1923 memoir on
affinely connected manifolds standing at the head of that line; the
axiomatic form in which the covariant derivative is now taught descends
from Koszul, and lives in the same standard references this series has
cited since Chapter 1.

One forward clause, so the reader knows a debt when they see one paid
later: among all carrying rules, once a metric arrives --- chapters from
now --- exactly one will be singled out as the unique rule that is
torsion-free and preserves the ruler's lengths under carrying; the trade
calls that the fundamental theorem of the subject, and this book will
collect it where the metric is bought, not before. Until then the
connection stays what this chapter built: chosen, not given; extra, not
implied; honest about being a decision.

And now the question that hands the next chapter its subject. A rule for
carrying invites an experiment no reader will be able to resist: choose a
closed route --- out from the bench, around, and home again --- and carry
an arrow all the way around it, faithfully, by the rule. The arrow comes
home. Does it come home \emph{as itself}? On the schoolroom's flat paper,
always. On the Earth of the equator-and-pole experiment, the reader
already knows the answer is no --- the arrow returns turned. The amount
of the turning, per loop, is not an accident of the route's clumsiness;
it is a measurement, it is linear in the loop's little spans, and the
machine that answers it --- the machine that eats a small closed loop and
reports what carrying around it does --- is the object this entire
construction has been walking toward. The next chapter builds it.
